\subsubsection{User stories}
Based on the analysis presented in the state of art section and the interview conducted with a user from the target group, a set of user stories has been developed. These user stories aim to clarify user needs and expectations and describe the functionalities that are essential for the two identified user roles: wearer and guardian.

\begin{enumerate}
    \item As a user, I want to press the emergency button so that I can alert my relatives. 
    \item As a user, I want to send an alert when I fall so that I can notify my relatives. 
    \item As a user, I want to be able to cancel an alarm so that I don't have an active alarm when I am not in an emergency.
    \item As a user, I want to be able to log in to the application so that I have an overview of my relatives and my alarms.
    \item As a user, I want to add multiple relatives to my notification list so that multiple relatives receive an alarm and help me. 
    \item As a user, I want to see the battery percentage in my app so that I can ensure my wristband is charged. 
    \item As a relative, I want to be notified when an alert has been triggered so that I can react to the alarm.
    \item As a relative, I want to know when, where and how alerts have been triggered so that I know the scope of the emergency.
    \item As a relative, I want to see the user's real-time \gls{gps} location so that I am aware of where the user is at all times.
    
\end{enumerate}

The user stories presented in this section capture the system’s functionality from the perspective of its primary actors and describe how different users are expected to interact with the system in practical scenarios. By focusing on user goals rather than technical implementation, the user stories help ensure that the system design remains aligned with real user needs and expectations.

These user stories form the foundation for identifying the core interactions between the actors and the system. In the following section, they are translated into requirements in combination with insights from the state of art section and presented in two tables: functional requirements and non-functional requirements.

